\section*{Методология}

Атрибутивный метаграф $G = (V, E, A_V, A_E)$ определяется как ориентированный граф, где:
\begin{enumerate}
    \item $V$ — множество метаузлов (объектов ИС: процессы, пользователи, хосты, файлы, сетевые соединения);
    \item $E \subseteq V \times V$ — множество дуг (событий: запуск, запись, передача данных, сетевое взаимодействие);
    \item $A_V: V \to \mathcal{A}_V$ — отображение узлов в пространство атрибутов (например, \texttt{user\_id}, \texttt{privilege\_level}, \texttt{file\_hash});
    \item $A_E: E \to \mathcal{A}_E$ — отображение дуг в атрибуты событий (время, параметры вызова, IP-адреса, данные payload).
\end{enumerate}

\vspace{0.5em}
Атрибутивный метаграф отличается от классического графа возможностью описывать как объекты, так и отношения между ними с использованием атрибутов, что делает его особенно подходящим для моделирования сложных сценариев компрометации. В отличие от простых графов угроз, метаграфы позволяют:
\begin{enumerate}
    \item представлять вложенные структуры (например, процесс содержит потоки, которые открывают сокеты);
    \item задавать логические условия на атрибуты («если \texttt{\seqsplit{process\_name = powershell.exe}} И \texttt{\seqsplit{parent\_process $\neq$ explorer.exe}}»);
    \item поддерживать динамическое обновление при поступлении новых событий.
\end{enumerate}

Каждой технике MITRE ATT\&CK (например, T1059 — \textit{\seqsplit{Command and Scripting Interpreter}}) сопоставляется \textit{паттерн метаграфа} — подграф с заданными атрибутивными ограничениями и топологией. Например, для T1059 паттерн включает:
\begin{enumerate}
    \item узел типа «процесс» с атрибутом \texttt{\seqsplit{name $\in$ \{cmd.exe, powershell.exe, wscript.exe\}}};
    \item входящую дугу от узла «пользователь» или «процесс»;
    \item отсутствие родительского процесса из доверенного списка (whitelist).
\end{enumerate}

Обнаружение атаки сводится к задаче \textit{поиска изоморфного вложения} паттерна в текущий метаграф системы. Для этого используется алгоритм VF2, модифицированный для учёта атрибутивных ограничений: совпадение узлов проверяется не только по топологии, но и по заданным порогам схожести атрибутов (например, хэши файлов, уровни привилегий).

Процесс реализуется в три этапа:
\begin{enumerate}
    \item \textbf{Сбор и нормализация событий} из источников (аудит ОС Windows/Linux через Sysmon или Auditd, EDR-агенты, сетевые логи Zeek, журналы приложений) → унифицированные события в формате, совместимом с SIEM (например, CEF или ECS).
    \item \textbf{Построение динамического метаграфа} в реальном времени. При поступлении нового события система либо добавляет новый узел/дугу, либо обновляет существующие атрибуты (например, при повторном запуске процесса). Метаграф хранится в памяти в виде индексированной структуры данных (например, на основе GraphBLAS).
    \item \textbf{Сопоставление с библиотекой паттернов}, основанной на ATT\&CK и внутренних угрозах. Поиск выполняется с временным окном (по умолчанию — 5 минут), чтобы учитывать последовательности событий в рамках одной атаки.
\end{enumerate}

Для минимизации ложных срабатываний введены \textbf{контекстные условия}: например, запуск PowerShell из веб-браузера считается подозрительным, тогда как из консоли администратора — нет. Также учитывается временной контекст: последовательность событий должна укладываться в временные рамки, заданные для конкретной тактики.

\begin{figure}[ht]
\centering
\resizebox{0.9\textwidth}{!}{\begin{tikzpicture}[
  node distance=1.8cm,
  metanode/.style={rectangle, draw=blue!70, fill=blue!5, rounded corners, minimum width=1.8cm, minimum height=0.7cm, font=\small},
  attr/.style={font=\scriptsize, color=gray},
  edge/.style={->, thick, >=Stealth}
]

\node[metanode] (U) {User};
\node[metanode, below right=of U] (P) {Process};
\node[metanode, below left=of P] (N) {Network};

\draw[edge] (U) -- (P) node[midway, right, attr] {executes};
\draw[edge] (P) -- (N) node[midway, below, attr] {connects};

\node[attr, above=0.1cm of U] {uid=admin};
\node[attr, right=0.1cm of P] {name=powershell.exe};
\node[attr, below=0.1cm of N] {dst\_ip=192.168.10.55};

\end{tikzpicture}}
\caption{Пример фрагмента атрибутивного метаграфа: пользователь → процесс → сетевое соединение}
\end{figure}